\subsection{System Architecture}
\label{sec:system}

Fig.~\ref{fig:system-overview} shows the stack.
\textbf{Device layer:} FR3 manipulator, RGB camera, and ROS~2/MoveIt PTP execution; named skills invoke calibrated joint poses via the edge gateway (\texttt{PAPER1\_MODE=1} disables stiffness routes in experiments).
\textbf{Edge layer:} FastAPI gateway (\texttt{franka\_api\_server}) hosts Edge-IQA and LangGraph; \texttt{sim/franka\_bridge} connects graph nodes to \texttt{/vision/evaluate} and skill endpoints.
\textbf{Cloud layer:} higher-level vision and EMR integration (stub); upload occurs only when edge quality and routing accept the frame.

The boundary between control and perception is explicit.
Edge-IQA never writes directly to the 1\,kHz hardware loop; it only selects among pre-calibrated skills or a fail-safe return path exposed by the gateway.
Thus a perception fault degrades the acquisition workflow, not the low-level controller.
This design mirrors RA-L systems that place the algorithmic contribution inside a real robot execution envelope while keeping the safety-critical path simple enough to audit.

Each trial logs ISO~8601 timestamps (\texttt{t\_capture}, \texttt{t\_iqa}, \texttt{t\_route}) and decisions to JSONL; main RTT metrics are computed offline by \texttt{run\_matrix\_tcm.py} on ShezhenV3 test images.
Offline RTT aggregates a documented latency model: capture 5\,ms, IQA $\mathcal{N}(9,2^2)$\,ms, routing 1.5\,ms, resample p50$\approx$50\,ms, cloud upload 50--550\,ms (see \texttt{sim/NETEM.md}).
Relative baseline ordering is the primary claim; ShezhenV3 blur labels are synthetically injected Gaussian defocus when human quality annotations are unavailable.
The same JSONL schema is used for offline replay and the auxiliary Franka track, so every uploaded, retried, or aborted frame can be traced to a score, flag set, retry count, and timestamp.

\begin{figure}[!t]
  \centering
  \includegraphics[width=\linewidth]{fig1_system_overview.pdf}
  \caption{Cloud--edge--device stack for robotic tongue-image acquisition.}
  \label{fig:system-overview}
\end{figure}
